\section{The canonical prime-carrying inverse has no bounded narrow-strip extension}
\label{sec:mobius-unbounded}

The finite-Euler remainder of MA-TR-101 rules out one proposed
limit argument. A different route gives a stronger obstruction
for the actual smooth theta--M\"obius inverse itself. It uses
small values of zeta on a fixed vertical line, not the location
of any zero. The obstruction therefore persists even when the
Dirichlet theorem $A$ is available.

\paragraph{MA-MO-001 (unboundedness of the canonical inverse).}
Let $0<\alpha<1/2$ and $s\ge0$, and give both the radial
half-density variable $F$ and the centered theta variable
$g=\mathcal ZF$ the norm
\[
 \|u\|_{\alpha,s}^2=
 \int_{\mathbb R}|u(x)|^2
             e^{2\alpha|x|}(1+x^2)^{s/2}\,dx,
 \qquad x=\log t.
 \tag{MA-MO-1}
\]
The inverse $\mathcal J$ in (MA-TR-7), on the actual smooth
theta image, has no bounded extension from
$J_{\alpha,s}$ to $H_{\alpha,s}$ agreeing with that inverse.
Equivalently, no constant $C$ satisfies
\[
        \|F\|_{\alpha,s}\le C\|\mathcal ZF\|_{\alpha,s}
 \tag{MA-MO-2}
\]
even on the smaller domain
$\{F\in\mathcal E:\mathcal MF(1/2)=0\}$.
Consequently there are smooth pairs $g_j=\mathcal ZF_j$ with
$\|F_j\|_{\alpha,s}=1$ and $\|g_j\|_{\alpha,s}\to0$.

\begin{proof}
First this smaller domain is contained in the one used in
MA-TR-101. If $F$ is exponential Schwartz as a function of
$x=\log t$, set $f(t)=F(\log t)/(2\sqrt t)$ for $t>0$,
extend it evenly, and set $f(0)=0$. Every derivative tends
to zero at the origin and decays faster than every power at
infinity; this follows by differentiating in $t$, which only
introduces fixed powers of $t^{-1}$ and logarithmic derivatives
of $F$. Thus $f$ is even Schwartz and flat at zero. Moreover
\[
       \int_{\mathbb R}f(t)\,dt=\mathcal MF(1/2)=0.
\]
MA-TR-101 therefore gives $g=\mathcal ZF\in\mathcal E$ and
the exact identity
\[
 \mathcal Mg(z)=\zeta(z+1/2)\mathcal MF(z).
 \tag{MA-MO-3}
\]
Initially this follows by absolute convergence for $\Re z>1/2$;
the zero moment removes the only zeta pole, so analytic
continuation gives it everywhere. No zero-free region is used.

Suppose (MA-MO-2) holds. Write $U_XF(x)=F(x-X)$. It preserves
the domain because its Mellin transform is
$e^{Xz}\mathcal MF(z)$, and $\mathcal ZU_X=U_X\mathcal Z$.
For each fixed $F\in\mathcal E$,
\[
 \frac{\|U_XF\|_{\alpha,s}^2}
      {e^{2\alpha X}(1+X^2)^{s/2}}
 \longrightarrow
       \int_{\mathbb R}|F(u)|^2e^{2\alpha u}\,du
       \qquad(X\to+\infty).
 \tag{MA-MO-4}
\]
Indeed the ratio of weights tends pointwise to $e^{2\alpha u}$.
For $X\ge1$ it is bounded by
$2^{s/2}e^{2\alpha|u|}(1+u^2)^{s/2}$, using
$|X+u|-X\le|u|$ and
$1+(X+u)^2\le2(1+X^2)(1+u^2)$.
This is an integrable dominating weight against $|F|^2$.
The same argument applies to $g\in\mathcal E$.
Apply (MA-MO-2) to the translated pair and pass to the limit:
\[
 \int|F(x)|^2e^{2\alpha x}\,dx
 \le C^2\int|g(x)|^2e^{2\alpha x}\,dx.
\]
Plancherel on the vertical Mellin line and (MA-MO-3) give
\[
 \int_{\mathbb R}|\mathcal MF(\alpha+it)|^2\,dt
 \le C^2\int_{\mathbb R}
     |\zeta(1/2+\alpha+it)|^2
                  |\mathcal MF(\alpha+it)|^2\,dt.
 \tag{MA-MO-5}
\]
The common Fourier normalization constant cancels.

Fix any $t_0\in\mathbb R$ and put, for $\delta>0$,
\[
 \phi_\delta(x)=e^{-\delta x^2-\alpha x-it_0x},
 \qquad F_\delta=(\partial_x+1/2)\phi_\delta.
\]
These are in $\mathcal E$, and integration by parts gives
$\mathcal MF_\delta(z)=(1/2-z)\mathcal M\phi_\delta(z)$.
In particular their moment at $1/2$ vanishes. On the required
line, their squared Mellin transform is a positive constant
times
\[
 \bigl((1/2-\alpha)^2+t^2\bigr)
                     e^{-(t-t_0)^2/(2\delta)}.
\]
After dividing by its integral, this is an approximate point
mass at $t_0$, because $1/2-\alpha>0$.
The classical fixed-strip polynomial growth bound for zeta,
together with the Gaussian tail, justifies this limit also
against the factor $|\zeta(1/2+\alpha+it)|^2$.
Thus (MA-MO-5) implies
\[
          |\zeta(1/2+\alpha+it_0)|\ge C^{-1}
                     \qquad\hbox{for every }t_0\in\mathbb R.
 \tag{MA-MO-6}
\]

This contradicts the classical universality theorem of Voronin
(1975), the unnumbered theorem in \S I, pp.~443--444.
For clarity its consequence needed here is only
\[
       \inf_{t\in\mathbb R}|\zeta(\sigma+it)|=0
                    \qquad(1/2<\sigma<1).
\]
It follows directly from the theorem as printed, without using
a later stronger formulation: choose
$|\sigma-3/4|<r<1/4$, approximate the nonzero constant
$\varepsilon$ on the disc centered at $3/4$ with radius $r$
by a vertical translate of zeta to error $\varepsilon$, and
evaluate at its real point $\sigma$.
This proves the contradiction. Finally the failure of every
bound (MA-MO-2) supplies pairs with
$\|F_j\|>j\|g_j\|$; rescale to obtain the asserted sequence.
\end{proof}

\paragraph{Corollary: the original Schwartz target does not admit a continuous extension either.}
The inverse, now read as $g\mapsto f$ in
$\mathscr S_{0,+}\subset\mathscr S(\mathbb R)$, has no continuous
extension $J_{\alpha,s}\to\mathscr S_{0,+}$ for
$0<\alpha<1/2$. Indeed the map
$f\mapsto F(t)=2\sqrt t\,f(t)$ is continuous from this Schwartz
subspace into $H_{\alpha,s}$. Near zero, evenness and $f(0)=0$
give $|f(t)|\le(t^2/2)\sup_{|u|\le1}|f''(u)|$; at infinity
any sufficiently high Schwartz seminorm suffices. Explicitly,
\[
 \|F\|_{\alpha,s}^2
 =4\int_0^\infty |f(t)|^2
       \max(t^{2\alpha},t^{-2\alpha})
                    (1+\log^2t)^{s/2}\,dt,
\]
and the displayed bounds make this integral a continuous
Schwartz seminorm estimate (the near-zero estimate even allows
$\alpha<5/2$). Composing the proposed continuous extension
with this embedding contradicts MA-MO-001. This does not
exclude individually regularized maps obtained by additional
spatial and Fourier cutoffs.

\paragraph{Precise consequence for the proposed transport.}
The theorem rules out a bounded Hilbert-space realization of
the \emph{specified} inverse with these two natural radial
weights. Since a bounded extension is determined by its values
on the dense smooth geometric image, there is no different
bounded extension that agrees with that same arithmetic map
there. This is not an assertion that every prime-corrected
intertwiner between every possible pair of completions is
impossible. In particular, it leaves open an unbounded closed
operator, a target norm transported through $\mathcal Z$, a
different scale of spaces, and a directly controlled nuclear
regularization. Nor does the theorem contradict MA-LIFT-002,
which proves trace class for the actual convolution lift
without introducing this inverse as a bounded operator.

The source's missing step is still the passage in Section~12.3,
p.~514, used in the proof on p.~520, between the global weighted
geometric operator and the specified local one. MA-TR-101
constructs the smooth map with the needed arithmetic content;
MA-MO-001 shows that declaring this map bounded on the narrow
spaces does not fill that step. The remaining problem is now
to construct and estimate an \emph{unbounded or differently
completed} trace transport, including its prime correction,
or prove that no such transport realizes the source's kernels.
The universality input concerns small nonzero values. It uses
neither $A$, $M$, nor a hypothesized zero outside the critical
line, and hence is not a circular zero-location argument.

\paragraph{Grothendieck placement.}
This is an analytic completion obstruction for the adelic
theta map on $C_{\mathbb Q}$, in the precise smooth-quotient
identification of GRH\_NIGHT MT-M24. It refines
Q-MAS-GEOMETRIC-NARROW and ALPOGE\_BRIDGES R-MAS-KERNEL by
excluding the canonical bounded Hilbert transport.
In cote~1, scans~079--080 and 111--112, as located by
BR-TR-01/02 and F-SM-11/12, Fredholm trace and determinant
operations require specified topologies and nuclearity.
A smooth topological isomorphism does not supply those
properties in a different Hilbert completion. The theorem
does not construct a sheaf, a motive, or an algebraic
correspondence, and there is no additional EGA/SGA algebraic
theorem identifying the two analytic norms. The corresponding
$K_0$ and $\lambda$-ring identities do not remove the missing
norm estimate.

\paragraph{Relation to the spectral-operator literature.}
The one-sided small-value obstruction is established background.
Herichi and Lapidus (2014), Theorem~3.11, p.~641, identify the
spectrum of $\zeta(\partial_c)$ with the closure of
$\zeta(c+i\mathbb R)$; their \S6, p.~652, records that this
closure is all of $\mathbb C$ for $1/2<c<1$ by universality.
Thus its bounded invertibility fails even when a vertical line
contains no zero. Their Theorem~3.12 concerns a different,
truncated quasi-invertibility condition, which is equivalent to RH.
MA-MO-001's specific deduction is the translation-limit transfer
of the small-value obstruction to the two-sided polynomially
weighted theta-image norm and its original Schwartz target.
No novelty is claimed for the underlying one-sided spectral
obstruction.

\paragraph{Primary reference and check boundary.}
Voronin, S. M. (1975). Theorem on the ``universality'' of the
Riemann zeta-function. \emph{Mathematics of the USSR-Izvestiya,
9}(3), 443--453.
\url{https://doi.org/10.1070/IM1975v009n03ABEH001485}.
The exact theorem statement and its two printed pages were
freshly inspected; the classical universality proof is cited
as background, not asserted to have been independently
re-proved here. The deduction above is new work of this run.

Herichi, H., \& Lapidus, M. L. (2014). Truncated infinitesimal
shifts, spectral operators and quantized universality of the
Riemann zeta function. \emph{Annales de la Facult\'e des Sciences
de Toulouse: Math\'ematiques, 23}(3), 621--664.
\url{https://doi.org/10.5802/afst.1419}.
The cited statements and the concluding universality application
were freshly inspected; their full proofs are carried background.
